% Method — Implementation Details
% (Formal definitions are in Section~\ref{sec:notation}; this section describes the iterative protocol.)

\section{Implementation}
\label{sec:method}

We implement GCE using an iterative protocol with probe recycling. Each round of erasure begins by training a fresh concept scorer on the current (possibly already partially erased) representations, computing per-sample gradients, and applying the erasure step (Eq.~\ref{eq:ambient_erase} for Ambient GCE, Eq.~\ref{eq:tangent_erase} for TanGCE). The tangent estimator is refit in every round to track the evolving manifold geometry as representations shift.

Two diagnostics verify tangent quality throughout the iterative process. \emph{Tangent Gradient Capture} (TGC) measures the fraction of the leakage gradient that lies within the estimated tangent space. \emph{Tangent Neighbor Residual} (TNR) measures the fraction of local secant vectors not explained by the tangent basis. High TGC and low TNR confirm that the tangent estimation remains geometrically sound.
